\documentclass{article}

\usepackage[style=english]{csquotes}
\usepackage[british]{babel}

%% References
\usepackage[
    bibstyle=ieee,
    citestyle=numeric-comp,
    backend=biber,
    maxcitenames=2,
]{biblatex}
\addbibresource{../library_biblatex.bib}

%% Acronyms
\usepackage[automake]{glossaries-extra}
\setabbreviationstyle[acronym]{long-short}
\makeglossaries

% Apparently the only way to prevent all the acronyms from generating spelling errors is to 
% disable ltex on the entire file.
% LTeX: enabled=false

\newabbreviation{ai}{AI}{Artificial Intelligence}
\newabbreviation{ar}{AR}{Augmented Reality}
\newabbreviation{asf}{ASF}{Amplitude Stereo Features}
\newabbreviation{cac}{CAC}{Cut And Clicks}
% \newabbreviation{cd}{CD}{Cochleagram Difference}
\newabbreviation{cd}{CD}{Compact Disc}
\newabbreviation{dft}{DFT}{Discrete Fourier Transform}
\newabbreviation{dirac}{DirAC}{Directional Audio Coding}
\newabbreviation{dl}{DL}{Deep Learning}
\newabbreviation[longplural=Directions of Arrival]{doa}{DOA}{Direction of Arrival}
\newabbreviation{dsp}{DSP}{Digital Signal Processing}
\newabbreviation{erb}{ERB}{Equivalent Rectangular Bandwidth}
\newabbreviation{fa}{FA}{Frequency-Azimuth}
\newabbreviation{faa}{FAA}{Frequency-Azimuth Analysis}
\newabbreviation{fft}{FFT}{Fast Fourier Transform}
\newabbreviation{fir}{FIR}{finite impulse response}
\newabbreviation{fmttm}{FMTTM}{Fly Me To The Moon}
\newabbreviation{genai}{Gen-AI}{Generative Artificial Intelligence}
\newabbreviation{iamf}{IAMF}{Immersive Audio Model and Formats}
\newabbreviation{daw}{DAW}{Digital Audio Workstation}
\newabbreviation{dj}{DJ}{Disc Jockey}

\newabbreviation{rq}{RQ}{research question}

\newabbreviation{iid}{IID}{Interaural Intensity Difference}
\newabbreviation{ipd}{IPD}{Interaural Phase Difference}
\newabbreviation{icpd}{ICPD}{Inter-channel Phase Difference}
\newabbreviation{ccpd}{CCPD}{Cross-channel Phase Difference}
\newabbreviation{ild}{ILD}{Interaural Level Difference}
\newabbreviation{ilr}{ILR}{Interaural Level Ratio}
\newabbreviation{cld}{CLD}{Channel Level Difference}
\newabbreviation{ccad}{CCAD}{Cross-channel Amplitude Difference}
\newabbreviation{icld}{ICLD}{Inter-channel Level Difference}
\newabbreviation{itd}{ITD}{Interaural Time Difference}
\newabbreviation{ictd}{ICTD}{Inter-channel Time Difference}
\newabbreviation{iacc}{IACC}{Interaural Cross Correlation}
\newabbreviation{iccc}{ICCC}{Inter-channel Cross Correlation}
\newabbreviation{icc}{ICC}{Inter-channel Correlation}
\newabbreviation{ic}{IC}{Inter-channel Coherence}
\newabbreviation{iac}{IAC}{Interaural Coherence}
\newabbreviation{itf}{ITF}{Inter-channel Transfer Function}
\newabbreviation{lev}{LEV}{Listener Envelopment}
\newabbreviation{dsd}{DSD}{Directional Loudness Distortion}
\newabbreviation{bms}{BMS}{Bivariate Mixture Space}
\newabbreviation{spl}{SPL}{Sound Pressure Level}
\newabbreviation{iir}{IIR}{Infinite Impulse Response}


\newabbreviation{ishst}{ISHST\textsuperscript{**}}{I Saw Her Standing There\textsuperscript{**}}
\newabbreviation{kks}{KKS}{The Knowledge Foundation}
\newabbreviation{lb}{LB}{Laser Beam}
\newabbreviation{lfe}{LFE}{Low Frequency Effects}
\newabbreviation{llm}{LLM}{Large Language Model}
\newabbreviation{lufs}{LUFS}{Loudness Units Full Scale}
\newabbreviation{mirex}{MIREX}{Music Information Retrieval Evaluation eXchange}
\newabbreviation{mir}{MIR}{Music Information Retrieval}
\newabbreviation{ml}{ML}{Machine Learning}
\newabbreviation{mtg}{MTG}{Music Technology Group}
\newabbreviation{m}{M}{Manifesto}
\newabbreviation{pc}{PC}{Panning Coefficients}
\newabbreviation{ph}{PH}{Panning Histogram}
\newabbreviation{pi}{PI}{Panning Index}
\newabbreviation{picoc}{PICOC}{Population, Intervention, Comparison, Outcome, Context}
\newabbreviation{psf}{PSF}{Phase Stereo Features}
\newabbreviation{qm}{QM}{Qualitative Model}
\newabbreviation{rms}{RMS}{Root Mean Squared}
\newabbreviation{rm}{RM\textsuperscript{*}}{'Round Midnight\textsuperscript{*}} 
\newabbreviation{simpl}{SIMPL}{stereo-input mixed-to-panned level}
\newabbreviation{sip}{SIP}{Stereo Image Profile}
\newabbreviation{slr}{SLR}{Systematic Literature Review}
\newabbreviation{spf}{SPF}{Stereo Panning Features}
\newabbreviation{sps}{SPS}{Stereo Panning Spectrum}
\newabbreviation{ssps}{SSPS}{Spectral Stereo Phase Spread}
\newabbreviation{stft}{STFT}{Short-time Fourier Transform}
\newabbreviation{tf}{TF}{time-frequency}
\newabbreviation{titt}{TITT}{This Is The Thing}
\newabbreviation{vr}{VR}{Virtual Reality}
\newabbreviation{peaq}{PEAQ}{Objective Measurements of Perceived Audio Quality}
\newabbreviation{peaqmc}{PEAQ-MC}{Objective Measurements of Perceived Audio Quality - Multi Channel}
\newabbreviation{cnn}{CNN}{Convolutional Neural Network}

\newabbreviation{ica}{ICA}{Independent Component Analysis}
\newabbreviation{pca}{PCA}{Principal Component Analysis}
\newabbreviation{pqmf}{PQMF}{Pseudo Quadrature Mirror Filter}
\newabbreviation{mpeg}{MPEG}{Moving Picture Experts Group}
\newabbreviation{3gpp}{3GPP}{3rd Generation Partnership Project}

\newabbreviation{asw}{ASW}{Apparent Source Width}
\newabbreviation{fwhm}{FWHM}{Full Width Half Maximum}
\newabbreviation{lef}{LEF}{Lateral Energy Fraction}
\newabbreviation{iqr}{IQR}{Interquartile Range}
\newabbreviation{amt}{AMT}{Auditory Modelling Toolbox}
\newabbreviation{qestral}{QESTRAL}{Quality Evaluation of Spatial Transmission and Reproduction using an Artificial Listener}
\newabbreviation{bcc}{BCC}{Binaural Cue Coding}
\newabbreviation{ps}{PS}{Parametric Stereo}
\newabbreviation{saoc}{SAOC}{Spatial Audio Object Coding}
\newabbreviation{gccphat}{GCC-PHAT}{Generalized Cross Correlation with Phase Transform}
\newabbreviation{mp3}{MP3}{MPEG layer 3}
\newabbreviation{aac}{AAC}{Advanced Audio Coding}
\newabbreviation{ivas}{IVAS}{Immersive Voice and Audio Services}
\newabbreviation{sofa}{SOFA}{Spatially Oriented Format for Acoustics}

\newabbreviation{ir}{IR}{Impulse Response}
\newabbreviation{hrir}{HRIR}{Head-Related Impulse Response}
\newabbreviation{hrtf}{HRTF}{Head-Related Transfer Function}
\newabbreviation{midi}{MIDI}{Musical Instrument Digital Interface}
\newabbreviation{vi}{VI}{Virtual Instrument}
\newabbreviation{mushra}{MUSHRA}{MUlti Stimulus Hidden Reference and Anchor}
\newabbreviation{bsseval}{BSS\_EVAL}{Blind Source Separation Evaluation}
\newabbreviation{psd}{PSS}{Perceptual Spectral Difference}
\newabbreviation{bamq}{BAM-Q}{Binaural-Auditory-Model-based Quality}
\newabbreviation{tff}{TFF}{Time-Frequency Filter}
\newabbreviation{pco}{PCO}{Population, Context, Outcome}
\newabbreviation{wos}{WOS}{Web Of Science}
\newabbreviation{ss}{SS}{Semantic Scholar}
\newabbreviation{cpu}{CPU}{Central Processing Unit}
\newabbreviation{flop}{FLOP}{Floating Point Operation}

\newabbreviation{csv}{CSV}{Comma Separated Variable}
\newabbreviation{tsv}{TSV}{Tab Separated Variable}
\newabbreviation{api}{API}{Application Programming Interface}
\newabbreviation{doi}{DOI}{Digital Object Identifier}
\newabbreviation{url}{URL}{Universal Resource Locator}
\newabbreviation{prisma}{PRISMA}{Preferred Reporting Items for Systematic reviews and Meta-Analyses}
\newabbreviation{adress}{ADRess}{Azimuth Discrimination and Resynthesis}
\newabbreviation{agr}{AGR}{Automatic Genre Recognition}
\newabbreviation{som}{SOM}{Self Organising Map}
\newabbreviation{iresap}{IRESAP}{Information Retrieval in Embedded Systems for Audiovisual Artistic Processes}
\newabbreviation{mdu}{MDU}{Mälardalen Universitet}
\newabbreviation{kmh}{KMH}{Kungliga Musikhögskolan}




% datasets
\newabbreviation{boma}{BOMA}{{Bang \& Olufsen Music for Archimedes}}
\newabbreviation{uiowamis}{UIowaMIS}{University of Iowa Musical Instrument Samples}
% \newabbreviation{}{}{Bach10}
\newabbreviation{btpab}{BTPAB}{ByTeDance paired ambisonic and binaural}
\newabbreviation{dsd100}{DSD100}{Demixing Secret Database}
\newabbreviation{asascbrm}{ASASCBRM}{Database for Automatic Spatial Audio Scene Classification in Binaural Recordings of Music}
\newabbreviation{daoc}{DAOC}{Denon anechoic orchestral collection}
% \newabbreviation{}{}{FAIR-Play}
\newabbreviation{fma}{FMA}{Free Music Archive}
% \newabbreviation{}{}{MedleyDB}
% \newabbreviation{}{}{MedleyDB 2.0}
% \newabbreviation{}{}{MUSDB18}
% \newabbreviation{}{}{MUSDB18-HQ}
\newabbreviation{quasi}{QUASI}{QUaero Audio SIgnals}
\newabbreviation{rwc}{RWC}{Real World Computing}
\newabbreviation{sqam}{SQAM}{EBU Sound Quality Assessment Material}
% \newabbreviation{}{}{Slakh}
\newabbreviation{urmp}{URMP}{University of Rochester Multi-Modal Music Performance}
% \newabbreviation{}{}{ccMixter}
\newabbreviation{arsso}{ARSSO}{Anechoic Recording System for Symphony Orchestra}
\newabbreviation{rioossr}{RIOOSSR}{Recordings of Italian opera orchestra and soloists in a silent room}
\newabbreviation{msss}{MSSS}{Mixing Secrets for The Small Studio}
\newabbreviation{tss}{TSS}{Telefunken Studio Sessions}
\newabbreviation{mtgj}{MTGJ}{Music Technology Group Jamendo}
% \newabbreviation{}{}{Audio and Acoustic Signal Processing Technical Committee Online Resources}
\newabbreviation{maestro}{MAESTRO}{MIDI and Audio Edited for Synchronous TRacks and Organization}
% \newabbreviation{}{}{FiloSax}
% \newabbreviation{}{}{Gamelan Music Dataset}
% \newabbreviation{}{}{mobygratis}


% Journals
\newabbreviation{jaes}{JAES}{Journal of the Audio Engineering Society}
\newabbreviation{eurasip}{EURASIP}{European Association for Signal Processing}


\title{The ethics of perceptual listening tests}
\author{Richard Mitic}

\begin{document}

\maketitle

In the world of audio signal research, perceptual listening tests are among the most prevalent methods to validate a scientific claim. Their use dates back to very early research into psychoacoustics, for example the work of \citeauthor{jeffressPlaceTheorySound1948} who studied the mechanisms by which humans can perceive the directionality of sound \cite{jeffressPlaceTheorySound1948}. A typical perceptual test can take several forms (although some standards have emerged over time \cite{itu-r-bs-1534_MUSHRA,itu-r-bs.1116-1MethodsSubjectiveAssessment1997}) but the underlying mechanism is that a human subject will listen to a particular sound using specific audio reproduction equipment, and will subsequently answer a question about the sound. A frequent trope is to compare two or more versions of a sound and then give a preference as to which was the most pleasant. Another is to compare two sounds and answer simply \enquote{can you hear a difference between them?}. In general though, the subject can be asked anything about the audio material and their answer will be taken as a legitimate data point.

As listening tests deal with human subjects, it is naturally necessary to apply the same ethical principles as any other research field that deals with human subjects. As a high-level reference, one can refer to The Lean Research Framework \cite{hoffeckerLeanResearchFrameworkAgust2015}, which is focused on human-centric research. The framework outlines 4 principles that need to be applied to perceptual listening tests - \enquote{Rigorous, Respectful, Relevant, and Right-sized}. These core principles will form the foundations on which we assess the ethics of listening tests in this manuscript.

Let us first address the exact form of the tests themselves. As previously mentioned, several standard tests have arisen over the years. Notable examples are those which have been codified into international standards such as \cite{itu-r-bs-1534_MUSHRA} which measures \enquote{intermediate audio quality} and \cite{itu-r-bs.1116-1MethodsSubjectiveAssessment1997} which measures \enquote{small impairments} in audio quality. Not only do the standards explain the test methodology and format of the subjective questions, but also provide standards for the listening environment and audio reproduction equipment. As \citeauthor{mustajokiNewApproachResearch2017} states in \cite{mustajokiNewApproachResearch2017}, the importance of standardization in research is ever-present due to the increasing number of researchers and their ease of global mobility. Moreover, the way to achieve standardization is through \enquote{explicit and transparent dialogue}; One could argue that there is no dialogue more explicit and transparent than an international standard written by committee. This is not to say that other test methods are invalid, however. In fact, it is trivial to find examples of perceptual listening tests in academic literature that are described as \enquote{MUSHRA-like}, implying that the standards can be used as a customizable base, safe in the knowledge that other researchers will understand the procedures and therefore trust that the results are legitimate.

The design of a listening test must also involve the definition of what data should be collected. According to the Lean \enquote{Right-size} principle, only the data required to answer the research question at hand should be collected. In the context of audio, this typically means that data which could be classed as \enquote{personal} under GDPR \cite{DataProtectionExplained} should not be collected, if possible. Generally speaking, it is possible to design listening tests that do this since the subjects' opinions, not their personal information, are of interest. However, the age of the subject is usually an important data point since hearing deteriorates with increased age. While age alone cannot identify a person, it could be combined with complementary information in order to identify a person \cite{akermanGoodResearchPractice}. While it is not a hard-and-fast rule that no personal data should be collected, the experimental design of the listening test can easily be formed such that collecting any such data is not necessary. 

A basic ethical principle which underpins respect for the subject is the prevention of harm \cite{tauginieneGeneralGuidelinesAcademic2018}. During listening tests, an obvious risk to the subject is that of hearing damage caused by the reproduction equipment. In simple terms, this is mitigated by ensuring that the sound pressure level at the subject's ears does not exceed a dangerous level. The \gls{mushra} standard \cite{itu-r-bs-1534_MUSHRA} offers a simple technique to achieve this - give the subject a volume control, empowering them to ensure their own comfort. While it may seem obvious, it is also imperative to use commercial reproduction equipment that has been approved for general use according to local standards. For example, all headphones sold in the \gls{eu} require a \enquote{CE mark}, denoting that they comply with the relevant standards for consumer safety \cite{CEMarkingObtaining,LegislationRegulationProduct}. Similarly, in the United States all commercial electronics that deal with electromagnetic radiation (which includes headphones) are required to be approved by the \gls{fcc}. 

Once the experimental design is established, communication of the aforementioned concerns with the test subjects is essential to ensure that their consent is informed. \citeauthor{akermanGoodResearchPractice}'s document \citetitle{akermanGoodResearchPractice} \cite{akermanGoodResearchPractice} gives a checklist that consent criteria shall fulfil in order to satisfy ethical concerns. Two of the points are that consent is voluntary and possible to withdraw. For listening tests these are especially important because the subjects may experience \enquote{listening fatigue}. This refers to the concept that prolonged exposure to audio stimuli can induce tiredness in the subject due to the cognitive load \cite{ronnbergListeningEffortFatigue2014}. Not only does this constitute a risk of harm that should be mitigated, it also makes the subjects' responses less reliable. A general rule of thumb to avoid listening fatigue is that the test should last less than 30 minutes, although it is perfectly possible that a subject will experience fatigue sooner. In particular, it has been shown that subjects with hearing impairments will tire quicker than normal-hearing subjects \cite{ListeningFatigue}. The communication of the subject's right to withdraw from the test at any point is therefore paramount, as extending the test beyond the fatigue point benefits no-one.

In the case that a listening test terminates early, the subject then has a choice whether to withdraw the data that has already been collected, or not. This opens the possibility that the collected data may be incomplete. When a researcher is post-processing the data, there is therefore a need to account for missing data points. The exact method to do this depends entirely on the task at hand, but it should also be explicitly stated in the communication of the research that data was incomplete. In particular, attempting to extrapolate missing data may count as fabrication or falsification as defined by the ALLEA code of conduct \cite{allea-alleuropeanacademiesEuropeanCodeConduct2023}, resulting in a case of scientific misconduct. 

To satisfy the condition of \enquote{Rigour} \cite{hoffeckerLeanResearchFrameworkAgust2015}, the collected data requires appropriate statistical analysis. The \gls{mushra} standard \cite{itu-r-bs-1534_MUSHRA} gives guidelines for how to collect a dataset which allows such analysis. Firstly, it specifies that roughly 20 subjects is usually enough, although it may be possible to simply continue adding more subjects until a desired level of statistical significance is reached. Importantly for listening tests, it also defines a process for \enquote{post-screening} of the subjects, eliminating those who gave a high rating to a \enquote{significantly impaired anchor signal}. Likewise, it also requires a \enquote{hidden reference} signal that is not degraded, allowing the post-screening to eliminate subjects that rated it badly. This further emphasizes the importance of the experiment design in that it must be possible to include the anchor and reference signals in the comparisons.

The \gls{mushra} standard \cite{itu-r-bs-1534_MUSHRA} also gives guidance as to who the test subjects should be, stating that they should \enquote{have normal hearing} and \enquote{have experience in listening to sound in a critical way}. However, with the increasing globalization of research mentioned by \citeauthor{mustajokiNewApproachResearch2017} \cite{mustajokiNewApproachResearch2017}, the results of traditional audio tests are being questioned. Recent work by \citeauthor{chauhanCrossCulturalBiasMelScale2026} showed that in a task relying on \textit{objective metrics} where the audio material was sourced globally, there were implicit geographical biases encoded in the choice of audio processing frontend \cite{chauhanCrossCulturalBiasMelScale2026}. It follows that there is a possibility of similar bias when selecting listening test subjects. For example an \enquote{experienced critical listener} might imply that the person comes from a specific background or was educated in a particular system, in which case people from other all other systems are not represented by the collected data. This issue warrants further investigation and has not been discussed widely in the literature.

As is the case with much current research, it is common for audio research to be a co-production with industry. As an example, it is perfectly normal for a company selling lossy audio compression software to fund research into new compression algorithms (which will require listening tests to validate). According to the lean principles \cite{hoffeckerLeanResearchFrameworkAgust2015} the data collected by a listening test should be relevant to all stakeholders, in which case the specific questions asked of the subjects in a listening test need to be designed by all parties collaboratively and should be equally beneficial to all stakeholders \cite{pageEthicsCoproductionKnowledge2022}. Use of the ITU standard tests and is a good starting point for this collaborative design since they are well \enquote{battle-tested}. Nevertheless, this presents another issue that needs to be handled in the experimental design before the listening test takes place.

As is typical with collaborative research, listening tests will typically involve many researchers. As previously mentioned, not only is the experimental design is complicated, requiring input from all involved parties, but the collection of data is a slow and costly process that can benefit from multiple researchers being involved \cite{delgadoDesignChoicesBinaural2023}. Statistical analysis of the results presents and additional skill set and could be delegated to other researchers. All researchers involved will therefore be required to write and approve of the final article. In short, collaboration is almost always essential for a good listening test. Under the guidelines given by \gls{icmje} \cite{icmje_recommendations}, all researchers in this hypothetical situation would qualify as authors, and therefore the discussion of authorship needs to happen early in the planning process.

In conclusion, the ethical considerations of perceptual listening tests must begin from the very outset of the experimental planning, and the design of the audio test is highly influenced by them. Since they are human tests like any other, the standard conventions and regulations must be followed. Although there are specific consideration for listening tests, such as ensuring that there is no chance of hearing damage or listening fatigue.

\printglossaries

\printbibliography

\clearpage

\end{document}